\section{Summary and Future Directions}
\label{sec:summary}

%%% =========================================================================
%%% SECTION 6 OVERVIEW & NAVIGATION GUIDE
%%% Total Recommended Section Word Count: 380 -- 450 words
%%% Word budget allocation:
%%%   - Section 6.1 (Summary of Findings & Critical Discussion): 140 -- 160 words
%%%   - Section 6.2 (Post-Training with RLVR): 120 -- 140 words
%%%   - Section 6.3 (Mixture of Experts / MoE Exploration): 120 -- 140 words
%%% =========================================================================


% ---------------------------------------------------------------------------
% Subsection 6.1
% ---------------------------------------------------------------------------
\subsection{Summary of Empirical Findings and Limitations}
\label{sec:summary_findings}

The empirical evaluation demonstrates that causal transformers can learn chess grammar directly from coordinate sequences without hardcoded board tensors. Through 20 training epochs, Nebium-Small (117M) converged to a validation loss of 2.727 nats (perplexity 15.30) with 100\% unconstrained opening legality. Scaling to Nebium-Medium (345M) lowered validation loss to 2.376 nats (perplexity 10.76) and raised top-5 move accuracy to 70.5\%. Under dual-T4 compute and session-time limits, Nebium-Large (762M) was trained for 3 epochs, reaching 80.0\% move legality and 6.071 nats validation loss, tracking the early trajectory of the smaller variants at matched depth.

Two primary limitations persist. First, unconstrained decoding exhibits spatial coordinate drift over deep horizons ($>60$ plies), where self-attention fails to preserve piece occupancy without explicit board tracking. Second, single-pass forward inference cannot replace lookahead tree simulation during sharp tactical complications, where finding optimal lines requires multi-ply search rather than greedy probability sampling.


% ---------------------------------------------------------------------------
% Subsection 6.2
% ---------------------------------------------------------------------------
\subsection{Reinforcement Learning with Verifiable Rewards (RLVR) Post-Training}
\label{sec:future_rlvr}

Self-supervised pre-training aligns predictions with human game frequency rather than objective move quality. To address tactical errors, post-training via Reinforcement Learning with Verifiable Rewards (RLVR) presents a direct next stage. Unlike natural language, chess provides deterministic ground-truth verification through programmatic state engines: \texttt{python-chess} evaluates move legality instantly, while Stockfish provides centipawn evaluation deltas ($\Delta\text{CP}$) and terminal game outcomes.

This setting allows rule-based policy optimization using Group Relative Policy Optimization (GRPO) or Proximal Policy Optimization (PPO) without subjective human preference models. By sampling trajectory rollouts and assigning negative rewards to illegal moves or blunders exceeding 200 centipawns, policy gradients shift unconstrained generation toward legal, tactical accuracy. Such verifiable reinforcement learning targets coordinate drift in deep horizons while retaining the fast sampling speed of feedforward inference.


% ---------------------------------------------------------------------------
% Subsection 6.3
% ---------------------------------------------------------------------------
\subsection{Mixture of Experts (MoE) Architecture as an Empirical Test Bed}
\label{sec:future_moe}

Autoregressive chess modeling provides a controlled, mathematically bounded domain for investigating Sparse Mixture of Experts (MoE) architectures. By replacing the dense SwiGLU feed-forward sublayers with top-$k$ gated routing over 8 experts (activating 2 per token), total parameter capacity can scale beyond one billion weights while holding inference FLOPs matched to Nebium-Medium.

This environment enables empirical testing of routing specialization across distinct game phases. Because chess games shift between opening book theory, complex middlegames, and simplified endgames, gating routers can be inspected to determine whether individual experts specialize by phase. Quantifying router entropy across game plies and measuring expert activation frequency by piece type will reveal whether sparse conditional computation discovers strategic hierarchies without explicit phase annotations.
